% !TeX root = ../papers/nc-p3.tex
\section{全四隅の五行列恒等式と正則窓の三十二次判定}
\label{sec:five-matrix-finite-windows}
\begin{lemma}\label{lem:four-balanced-one-direction}
四辺有限な標準正則窓で頂点次元が全て4の単純 $S$ を取る。各頂点単純との第一Ext次元は
両向きとも3以下であり、全頂点から来る側と全頂点へ向かう側の少なくとも一方は零である。
行列恒等式は仮定しない。
\end{lemma}
\begin{proof}
関係空間 $R\subset V\otimes U$ の $V$ 支持を次元 $1,2,3$ の部分空間へ制限すると、
そこに含まれる関係の次元はそれぞれ $0,1,3$ 以下である。
純関係はなく、三次元支持の上界は一点評価階数3以上から従う。
二次元支持に二本あればその直交直線で残る式は四本以下となる。
有限辺の最大小行列式イデアルはその直線上では高々一本の四次式であるから、
直線が全て入るか長さ4の切断がある。いずれも有限Frobenius窓の直線切断長上界3に反する。

続く行の拡張の解空間を求める。直前の矢を各入力ベクトル $s$ で評価した
四行四列 $A_s:V\to k^4$ の四次元空間を取る（双対の規約を揃えて作用を読む）。
非零入力での階数は2以上、全像は $k^4$ を張る。解空間は
$\sum_s(A_s\otimes1)R$ の零化空間である。この和の次元が9以上であることを示す。
一つが可逆なら、鉛筆中の非零な特異行列も一つ取る。階数3ならその関係像は六次元で、
可逆な方の像におけるその三次元支持の部分は高々三次元。
階数2なら関係像は五次元以上で、共通部分は高々一次元である。
可逆な方の六次元と足していずれも9以上となる。
一般階数3で可逆元がなければ、全像が張ることから異なる像超平面を持つ二つを取れる。
双方の関係像は六次元で、それらの共通部分の支持二平面は、元の写像で引き戻すと
三次元である。共通部分は高々三次元となり所望の下界を得る。

残る四次元の定階数2行列空間は存在しない。$\mathbb P^3$ 上で
$\mathcal O^4\to\mathcal O(1)^4$ の像束を $E$ とすると、
$E,E^*(1)$ はともに大域生成される。従ってその第一Chern数は $0,1,2$。
0なら $E$ は自明で（独立な二切断の行列式が非零定数）、階数2の余核束の第三Chern類が
$4H^3$ となり矛盾する。2は転置で同じ。1なら階数2の核束の第三Chern類が
$(-1+2c_2(E))H^3$ となり消えない。ここで $c_2(E)$ は数で書いた。
よって行の第一Extは、十六次元から少なくとも九条件と元の四行を引いて高々3。
列側は反対環で同じである。均衡型では一般倍率間の第一・第二Extが消え、各頂点で $a_i b_i=0$。
両方向があれば定理\ref{thm:one-sided-balanced-small-simples}の最短区間の
両行列の非単射条件を適用する。今得た上界3で二次元・三次元の証明と同じ余次元評価になる。
生じる一般二倍率間の三段の非零積は補題\ref{lem:growth-fixed-middle}に反する。
\end{proof}

\begin{theorem}\label{thm:five-matrix-finite-windows}
四辺有限な標準有限Frobenius候補について、標準Koszul左右AS正則で各頂点隅が
五行五列行列の全恒等式を満たすことは、定理\ref{thm:four-matrix-all-regular}の
八次中心元による32次までの検査と一般隅の指数検査に同値である。
全ての該当窓は中心有限な合成整域で、一般隅の指数は2または4である。
\end{theorem}
\begin{proof}[残り得る単純の大きさ]
五行列条件と正則性を仮定する。全非頂点単純は均衡である。
どこかで5に達すれば五行列の非恒等中心多項式をその隅に評価して正次数中心元を
全体へ延ばせる（命題\ref{prop:all-window-corner-center}）。
従って全頂点と両側で全Extが零となり、Euler式で均衡する。
最小次元が3以下なら定理\ref{thm:generic-corner-small-simples}、
その他は全位置が4である。頂点と両側で全Extが零の単純をこの証明では直交する単純と呼ぶ。
全ての非頂点単純の第一Extは全体で高々片向きである。大きさ4には直前の補題を使う。
両向きの第一Extが零ならCY双対とEuler式で全Extが零となる。

大きさ $d$ の非直交単純と大きさ $c$ の直交単純で $c+d+1>5$ は
補題\ref{lem:periodic-tail-vertex-bridges}に反する。
薄い単純は常に直交するので、大きさ4の全単純も直交する。
3以上の直交単純があれば全単純が直交し、定理\ref{thm:pi-vertex-extension-criterion}を適用できる。
大きさ2の非直交単純の倍率軌道は有限個である。大きさ3の非直交単純の軌道も有限個であることを確認する。
補題\ref{lem:binary-simple-trace-escape}の曲線の極での同時倍率と格子を使う。
特殊化 $S_0$ の半単純型に非薄因子は高々一つで、その大きさは2か3。
薄い因子は直交し、この非薄因子も全体で高々片向きなので、
非頂点と頂点の混在は一般単純へ平滑化できない
（定理\ref{thm:main-regular-specialization}の厚みごとの非対角消去）。
全て頂点という場合は単元のtraceに反する。一般の $S$ からある頂点への第一Extが非零なら、
特殊化にもその向きが残るため、組成列のどれかで同じ向きが非零である。
それはただ一つの非薄因子で、逆向きは全因子で零である。
一方、選んだ薄い単純から $S$ への第一Extの非消滅は、同時倍率で源を頂点の和へ特殊化しても残る。
この矛盾は二次元の証明と同じで、逆向きも同様である。

直交する非薄単純がなければ全単純の倍率軌道は有限個になり、
定理\ref{thm:four-matrix-all-regular}の軌道商の冪零積とAS消滅の証明で矛盾する。
ゆえに中心有限でないとすれば、直交する大きさ2の単純が存在し、非直交の大きさ3は上の
$2+3+1$ の不等式で排除される。全非頂点単純は薄いか大きさ2となり、非直交の軌道は有限個である。
定理\ref{thm:one-sided-pi-noether-criterion}によれば、頂点への拡張を持つ $S_+$ と
頂点からの拡張を持つ $S_-$ がともに存在する。双方大きさ2で、その倍率軌道は交わらない。
\end{proof}

\begin{proof}[二次元以下のtrace環だけでは頂点射影を支えられない]
上の残余の場合に、次元3以上の既約族をなす直交二次元単純の指標があることを示す。
全ての斉次素商のtrace環が次元2以下と仮定する。
定理\ref{thm:coherent-pi-regular-domains}の素商抽出の部分を非零頂点射影に用いる。
ここで抽出にはcoherent性を使わない。斉次素商 $R$ のtrace拡大 $T=RC$ と中心分母
$xT\subset R$ を取り、$x$ に捩れがあればより小さい次元の素商を選び直す。
これで頂点射影内の非零部分加群 $N=mxT$ が得られ、$N$ は有限生成可換斉次trace環 $C$ 上有限である。
有限次元部分加群は右ASのsocle消滅に反する。

$\operatorname{Spec} C$ の原点外では各trace繊維の全単純が非頂点である。
一般の単純の各頂点次元は高々2だから、その特殊化の重みの中で
非頂点と頂点が混在すれば非頂点は薄いものだけであり、それらの直交性に反する。
全部頂点は非原点のtraceで排除される。さらに非直交の因子があればそれ一つが
重み1で繊維の被約traceを占める。従って非直交のあり得る非原点指標は
有限個の倍率軌道に限られる。被約traceは $C$ の全生成元の値を決めるためである。

全ての非零繊維で頂点単純の商がないので、$N/NB_{>0}$ は原点支持で有限次元であり、
次数付き中山の補題から $N$ は $B$ 上有限生成である。
$C$ の正次数斉次元 $z$ を、先の有限個の非直交軌道上で非零であり、
かつ $N$ の非零因子となるように選ぶ。原点を支持とする随伴素点は有限次元部分加群の不在で除かれ、
残りの有限個の条件は次数を公倍数に揃えて素回避できる。
$Q=N/zN$ の $C$ 台は次元1以下で、原点外の全繊維の組成因子は直交する。
従って $\operatorname{Tor}^B_j(Q,\bigoplus k_i)$ は全 $j$ で有限次元である。
実際左頂点Koszul分解とテンソルして有限 $C$ 加群の複体を取り、
非原点の各極大イデアルの冪で割ると、組成列で全Torが消える。
完備化への極限は各有限長商のMittag--Leffler条件で完全性を保つ。
よってそのホモロジーは原点支持である。$Q$ の極小次数付き分解は全項有限生成となる。

$Q$ の有限次元捩れを除くと、残る次元1以下の $C$ 加群に正次数非零因子が取れ、その剰余は有限次元である。
有限次元部分、剰余はいずれも $B$ 上有限完全で、$B$ への次数付きExtは第4未満で消える。
有限次元部分を除いた加群も有限完全なので、その第0,1 Extは内部次数で下に有界である。
正次数非零因子の完全列でこれらが非零次数移動と同型になるから零である。
有限次元部分を戻しても $\underline{\operatorname{Ext}}^i_B(Q,B)=0$ ($i=0,1$)。
$0\to N(-\deg z)\to N\to Q\to0$ により今度は
$\underline{\operatorname{Hom}}_B(N,B)$ が非零次数移動と同型になる。
これは $N$ が $B$ 上有限生成ゆえ下に有界で、包含写像により非零である。矛盾を得た。

従ってある素商は次元3以上のtrace環を持つ。その一般開では中心分母を逆にして
一般行列環のAzumaya軌跡を取れ、異なる閉点は異なるtraceの単純を与える。
薄い型と非直交型は有限個の倍率軌道なので、この三次元以上の族の一般点は直交二次元単純である。
\end{proof}

\begin{proof}[反対向きの二つの単純に同時に接続する中間単純]
上の直交単純のそれぞれを $D$ とする。補題\ref{lem:periodic-tail-vertex-bridges}の
尾部の固有値の証明により、ある倍率で $D\to S_+$ の非零第一Extがあり、
別の倍率で $S_-\to D$ の非零第一Extがある。後者を全体の倍率で移せば、次の構成可能集合を得る。
\[
 \{(t,[L]):t\in\mathbb G_m,\ L\text{ は直交二次元単純},       \operatorname{Ext}^1((S_-)_t,L)\ne0,       \operatorname{Ext}^1(L,S_+)\ne0\}.
\]
構成可能性は有限Koszul複体の行列階数による。その $L$ 像をもう一つの倍率で動かした集合は
先の三次元以上の族を含む。よって表示した集合の次元は少なくとも2であり、
ある固定した $t$ の繊維には無限個の相異なる $[L]$ がある。
任意の有限個を選び補題\ref{lem:pi-yoneda-bridges}のYoneda積に用いる。
三つの単純は各頂点で $2+2+2>5$ なので、これら全中間因子から来るテンソル積の直和は
$\operatorname{Ext}^2((S_-)_t,S_+)$ に単射となる。この固定空間は有限次元であり矛盾する。

以上で中心有限性を得た。一般隅の指数は5以下、1と3は既出の有限グラフの排除、
5は定理\ref{thm:odd-prime-finite-indices}で除かれる。指数2または4の中心有限整隅は
四行列の全恒等式を満たす。従って定理\ref{thm:four-matrix-all-regular}の有限判定に帰着する。
逆にその判定を通れば四行列の恒等式を満たすので五行列の恒等式も満たす。
\end{proof}
